
		
		We start this section on symmetric functions by introducing, 
		\emph{semistandard Young tableaus}. 
		The reason for this is that semistandard Young tableaus are in bijection with 
		Gelfand-Tsetlin patters, 
		as shown later in Section \ref{gtpattern}. 
		We then continue through generalizing, introducing \emph{Schur functions}, 
		and some characteristics of Schur functions. 
		The source for this part is \cite{StanleyEC2}.
		
		\subsection{Semistandard Young tableaus} \label{SSYTsubsection}
		
		\definition{
			A \emph{semistandard Young tableau} (SSYT) of \emph{shape} $\lambda$ is an array $T = (T_{ij})$ of positive integers 
			such that it is weakly increasing in every row and strictly increasing down in every column.
		}
		
		\definition{
			The \emph{size} of an SSYT is its number of entries.
		}
		
		\bigskip
		
		We give an example to make SSYTs clearer. 
		
		\example{
			Below is a SSYT of shape $(5,4,2,2)$: 
			
			\begin{center}
			\begin{tabular}{c c c c c}
				1 & 1 & 1 & 3 & 4 \\
				2 & 3 & 3 & 4 & \\
				3 & 4 & & & \\
				4 & 5. & & &
			\end{tabular}
			\end{center}
			
			As we can see it has 5 integers in the  first row, 4 integers in the second row and 2 in the third and fourth row. 
			It is also weakly increasing in the rows and strictly increasing in the columns. 
			
			Usually a SSYT is written as a \emph{Young diagram}, with boxes filled with positive integers as such
			
			\begin{center}
			\begin{ytableau}
				1 & 1 & 1 & 3 & 4 \\
				2 & 3 & 3 & 4 & \none \\
				3 & 4 & \none & \none & \none \\
				4 & 5 & \none & \none & \none
			\end{ytableau}
			\end{center}
			
		}
		
		\definition{
			The SSYT $T$ has the \emph{type} $\alpha = (\alpha_1, \alpha_2, ...)$ if $T$ has 
			$\alpha_i = \alpha_i(T)$ parts equal to $i$.
		}
		
		\example{
			The SSYT from earlier, 
			
			\begin{center}
			\begin{ytableau}
				1 & 1 & 1 & 3 & 4 \\
				2 & 3 & 3 & 4 & \none \\
				3 & 4 & \none & \none & \none \\
				4 & 5 & \none & \none & \none
			\end{ytableau}
			\end{center}
			
			has type $(3, 1, 4, 4, 1)$ since there are $3$ $1$'s, $1$ $2$'s etc.
		}
		
		\definition{
			Let $\mu \subseteq \lambda$ such that $\mu_i \leq \lambda_i$ for all $i$. 
			
			We define a semistandard Young tableau of \emph{skew} shape $\lambda / \mu$ to be 
			an array $T = (T_{ij}$ of positive integers of shape $\lambda / \mu$ that is weakly 
			increasing in every row and strictly increasing in every column.
		}
		What we mean by this definition is that the elements in the set $\mu$ are removed 
		from the semistandard Young tableau, while all other rules apply as usual. 
		We let an example demonstrate. 
		
		\example{
			Below is a SSYT of shape $(5,4,2,2)/(2,1)$: 
		
			\begin{center}
			\begin{ytableau}
				\none & \none & 1 & 3 & 4 \\
				\none & 3 & 3 & 4 & \none \\
				3 & 4 & \none & \none & \none \\
				4 & 5 & \none & \none & \none
			\end{ytableau}
			\end{center}
			
			As we can see the SSYT is still weakly increasing in the rows and strictly increasing 
			in the column. 
			However, the first row is missing two integers in the beginning, 
			and the second row is missing one integer in the beginning, since 
			$\mu = (2,1)$.
		}
		
	\subsection{Schur functions}
	\emph{
		Semistandard Young tableaus are a combinatorial object underlying Schur functions. 	}

		\definition{
			We denote a \emph{Schur function} as $s_{\lambda}$.
		}
		
		\bigskip		
		
		We continue by introducing some definitions, theorems and propositions for Schur functions, using SSYTs as examples. 
		
		\definition{
			Let $\lambda / \mu$ be a skew shape. 
			The \emph{skew Schur function} $s_{\lambda / \mu} = s_{\lambda / \mu}(x)$ of 
			\emph{shape} $\lambda / \mu$ in the variables $x = (x_1, x_2, ...)$ is the formal 
			power series
			\[
			s_{\lambda / \mu} (x) = \sum_T x^T,
			\]
			summer over all SSYTs $T$ of shape $\lambda / \mu$. 
			If $\mu = \emptyset$, we call $s_\lambda (x)$ the \emph{Schur function} of shape 
			$\lambda$.
		}
		
		
		\theorem{
			For any skew shape $\lambda / \mu$, the skew Schur function $s_{\lambda / \mu}$ is 
			a symmetric function.
		}
		
		\proof{
			This was shown by Cauchy in \cite{cauchy1815memoire}.
		}
		
		\example{
			All SSYTs of shape $(2,1)$ with the largest part being at most 3 are 
			
			\[
			\begin{array}{cccccccc}
			\begin{ytableau}
			1 & 1 \\
			2 & \none
			\end{ytableau}
			&
			\begin{ytableau}
			1 & 2 \\
			2 & \none
			\end{ytableau}
			&
			\begin{ytableau}
			1 & 1 \\
			3 & \none
			\end{ytableau}
			&
			\begin{ytableau}
			1 & 3 \\
			3 & \none
			\end{ytableau}
			&
			\begin{ytableau}
			2 & 2 \\
			3 & \none
			\end{ytableau}
			&
			\begin{ytableau}
			2 & 3 \\
			3 & \none
			\end{ytableau}
			&
			\begin{ytableau}
			1 & 2 \\
			3 & \none
			\end{ytableau}
			&
			\begin{ytableau}
			1 & 3 \\
			2 & \none
			\end{ytableau}
			\end{array},
			\]
			
		which gives the Schur function 
		\[
		s_{21}(x_1,x_2,x_3) = x_1^2 x_2 + x_1 x_2^2 + x_1^2 x_3 + x_1 x_3^2 + x_2^2 x_3 
		+ x_2 x_3^2 + 2 x_1 x_2 x_3.
		\]
		It is easy to see that $s_{21}$ is a symmetric function. 

		}
		
		\definition{
			Let $\lambda$ be a partition of some integer $n$ and $\alpha$ a weak composition of $n$ 
			(i.e. a sum of non-negative integers adding up to $n$). 
			The \emph{Kostka number} $K_{\lambda \alpha}$ then denotes the number of SSYTs of 
			shape $\lambda$ and type $\alpha$.
		}
		
		\definition{
		In a similar way we define the \emph{skew Kostka number} $K_{\lambda / \mu, \alpha}$ 
		as the number of SSYTs of shape $\lambda / \mu$ and type $\alpha$. 
		}
		
		\proposition{
		\emph{
			Suppose that $\mu$ and $\lambda$ are partitions with $|\mu| = |\lambda|$ (i.e. they 
			are the same size) and $K_{\lambda \mu} \neq 0$. Then $\mu \leq \lambda$, refered to as \emph{dominance order}. 
			Moreover, $K_{\lambda \lambda} = 1$.
		}
		}
		
		\begin{proof}
			It is obvious that there only exists one SSYT of type the skew shape $\lambda / \lambda$. 
			
			Consider the SSYT with the shape $(3,2,1)/(3,2,1)$. 
			This is a SSYT with no integers, and this can only be done in one way.
		\end{proof}
		
		\begin{example}
			Consider the SSYT with $\lambda = (3,2,2)$, $\mu = (1,1)$, and $\alpha = (2,2,1)$. 
			Then the Kostka number $K_{\lambda / \mu, \alpha} = 2$, since there are two possible 
			SSYTs. These are 
			\[
			\begin{array}{cc}
			\begin{ytableau}
			\none & 1 & 1 \\
			\none & 2 & \none \\
			2 & 3 & \none
			\end{ytableau}
			& \text{ and} \quad
			\begin{ytableau}
			\none & 1 & 2 \\
			\none & 2 & \none \\
			1 & 3 & \none
			\end{ytableau}
			\end{array},
			\]
			since they need to be strictly increasing in the columns 
			and then we only have the choice of switching place on a $1$ and a $2$. 
		\end{example}
	
	\subsection{Littlewood-Richardson coefficients}
	\emph{
		Littlewood-Richardson coefficients appear in combinatorics in the theory of symmetric functions \cite{Rassart2004}.
	}
		\definition{
			The \emph{Littlewood-Richardson coefficients} $c_{\lambda \mu}$ are the constants in the multiplication rule for Schur functions defined as 
			\begin{equation}
				s_{\lambda} \cdot s_{\mu} = \sum_v c_{\lambda \mu}^v s_v.
			\end{equation}
		}